\section{Quantitative one-variable estimates}\label{sec:07-scalar}

This section proves the scalar bounds used in the matrix argument and in
Appendix~\ref{sec:08-segments}. The proof separates three tasks: exact
arithmetic at finitely many nodes, analytic interpolation between those
nodes, and estimates on the two unbounded tails. The final pointwise
inequalities then follow from a finite collection of rectangles with
strict positive margins.

\subsection{Profiles and the required estimates}

Throughout this section and Appendix~\ref{sec:08-segments}, finite decimal constants (including decimal endpoints of intervals) denote exact rationals. All notation introduced for the calculations below is local to this section. The parameters are
\[
\begin{gathered}
 b=.602,\quad c=.365,\quad k=\sqrt{b-c},\quad r=.22,\quad s=3.6,\quad w=4.6,\\
 \eta=.022,\quad \kappa=1.16,\quad a_R=7.5,\quad \lambda=.65.
\end{gathered}
\]
Put \(t_-= -.022,\ t_+=.064,\ t_c=.021,\ t_r=.043,\ m_*=.352\), so \(t_\pm=t_c\pm t_r\). Write \(\phi,p\) for the standard Gaussian density and distribution function and \(\mathcal N=x\partial_x-\partial_x^2\). We recall the following profiles on the real line (these formulas also specify them at zero):
\[
\begin{aligned}
 X=p/\phi,\quad Y=(1-p)/\phi,\quad f=-(1-p)-\log p,\quad j=-p-\log(1-p),\quad B=1-xY,\\
 d=(1+xX)^2 f+B^2j,\qquad g=\tfrac12(1-X^2f-Y^2j),\qquad
 v=\log(XY),\\
 \mathsf K=d-2g,\quad \mathsf C=-d+(a_R-1)g'',\quad t=\operatorname{arsinh}(x/s),\quad a_t=s\cosh t,\\
 \pi(x)=r(\cos(wt),\sin(wt)),\qquad \ell(x)=\sqrt{b-r^2-d(x)},\quad q(x)=k-\ell(x),\\
 S_q=(2+\mathcal N)q,\quad S_\pi=(2+\mathcal N)\pi,\quad \rho=|S_\pi|.
\end{aligned}
\]
Thus \(v=-\log[(\phi/p)(\phi/(1-p))]\). In numerical constants such as \(\sqrt{2\pi}\), the symbol \(\pi\) denotes the circular constant; the vector profile is denoted by \(\pi(x)\). Dots denote \(t\)-derivatives, \(n_F=\ddot F-\tanh(t)\dot F=a_t^2 F''(x)\).


The square root defining \(q\) is initially understood wherever its
radicand is positive. We prove a uniform lower bound for that radicand
before using any global derivative estimate for \(q\). Reflection makes
\(d,g,v,\mathsf K,\mathsf C,q\) even; the two components of \(\pi(x)\)
are even and odd, respectively.

\begin{lemma}[Real profile bounds]\label{lem:scalar-ranges}
For every real \(x\), one has \(\ell>0\). Bounds on values, and absolute bounds on derivatives, are as follows (a dash means no bound asserted):
\[
\begin{array}{c|cc|ccc}
 &\inf\ \ge &\sup\ \le & |\dot F|& |\ddot F| & |n_F| \\ \hline
 \mathsf K&-.007&.0275&.066&.238&-\\
 \mathsf C&-.501&-.341&.247&.79&.79\\
 q&.0775&.2559&.193&.53&.53
\end{array}
\]
Also \(|v'(x)-x|\le .319\), and
\[
 \lambda t_-^2+t_-\mathsf C+\mathsf K>0,\qquad
 \lambda t_+^2+t_+\mathsf C+\mathsf K<0,\qquad
 2(-.37)\mathsf C-(.37)^2-4\lambda\mathsf K > m_*^2.
\]
We have
\[
 U_*:=1.507\max(0,-n_{\mathsf K})+
 \frac{12\cdot .80}{m_*^2}
 \left[1.25(\dot{\mathsf K}+t_c\dot{\mathsf C})^2+5t_r^2 \dot{\mathsf C}^2\right]
 \le .54\,r^2 w^2 .
\]
The following additional bounds hold on \(|t|\ge1.4\):
\[
\begin{aligned}
 -.0001\le\mathsf K\le.0191,\quad -.501\le\mathsf C\le-.495,\quad .235\le q,\qquad
 |\dot{\mathsf K}|\le .0211,\quad |\dot{\mathsf C}|\le .0084,\quad n_{\mathsf K}\ge0.
\end{aligned}
\]
For \(|t|\ge3.6\), \(0\le\mathsf K\le .0008\), \(|\mathsf C+.5|\le .0002\), \(q\ge .255\). For \(|x|\ge5\), \(q\ge.2203\). Also
\(\int_0^\infty (\mathsf K)_+(x)\,dx\le.422\) and \(\int_0^{s\sinh(1.4)}(\mathsf K)_+\le .126\).
\end{lemma}

\begin{lemma}[Pointwise layer inequalities]\label{lem:scalar-layer}
We have \(.113\le S_q\le .514,\ \rho\le1.104\). For any \(z,x\), put
\[
 H=4c+4kq(x)+2(k-q(x))S_q(z)-2S_\pi(z)\cdot\pi(x),
 \qquad J=H+v''(x)-2\kappa,
 \quad b_m=\frac{(1+t_+)(b-3\eta)}{3(3(b+\eta)-4c)}.
\]
Thus $J$ is the profile $H_0$ in Proposition~\ref{prop:scalar-input}.
With \(h_s=.5,\ e_0=.256,\ \alpha=2.26\),
\[
 H+t_cJ-\tfrac12 t_r(h_s+J^2/h_s)-J^2/\alpha >
 2(b+\eta+(b+\eta-\kappa)t_-)+2b_m S_q(z)S_q(x)+2e_0(2t_r)^2 .
\]
Also \(b+\eta+\kappa t_-+.125>0\), \(1.053\kappa^2/[8(b+\eta+\kappa t_-+.125)]<e_0\) and \(e_0+(\log2+.25)\alpha/4<.80\). We also have
\(4c+2k(S_q(x)+S_q(z))-S_q(x)S_q(z)-S_\pi(x)\cdot S_\pi(z)>.25\).

\end{lemma}

We prove Lemmas~\ref{lem:scalar-ranges} and~\ref{lem:scalar-layer}
in the remainder of the section. In particular, none of their numerical
conclusions is used to justify the finite arithmetic that establishes it.

\subsection{Stable formulas and finite certificates}

By reflection it suffices to calculate on \(x\ge0\). Write \(u=1-p=Y\phi\),
\(h=h(u)=(-\log(1-u)-u)/u^2\), with the removable value \(h(0)=1/2\). The following formulas avoid cancellation between exponentially large factors:
\[
\begin{array}{ll}
 Y=\displaystyle\int_0^\infty e^{-xy-y^2/2}dy,\quad B=1-xY,& j=x^2/2+\log\sqrt{2\pi}-\log Y-p,\\
 d=B^2 j+(1-pB)^2 h,& \mathsf K=d+Y^2(j+p^2h)-1,\\
 g'=Y(Bj-(1-pB)p h),& g''=xg'-\mathsf K,\\
 R_v=v'-x=\phi/p-B/Y,& V_2=v''=2-(\phi/p)(x+\phi/p)-B/Y^2,\\
 D_1=d'=x(2d-1)-R_v-2g',&D_2=d''=2d+2xD_1-2g''-V_2,\\
 G_3=g'''=xg''+3g'-D_1,&G_4=g''''=xG_3+4g''-D_2 .
\end{array}
\]
Indeed \(X'=1+xX,\ Y'=xY-1,\ f'=-\phi(1-p)/p,\ j'=\phi p/(1-p)\); differentiation gives the identities. For \(F=d,\mathsf K,\mathsf C\), respectively, use
\[
\begin{gathered}
 (P_F,Q_F)=(D_1,D_2),\quad (D_1-2g',D_2-2g''),\quad
 (-D_1+6.5G_3,-D_2+6.5G_4),\\
 \dot F=a_t P_F,\qquad \ddot F=a_t^2 Q_F+xP_F,\qquad n_F=a_t^2 Q_F.
\end{gathered}
\]
Further,
\[
\begin{gathered}
 \dot q=\frac{\dot d}{2\ell},\qquad
 \ddot q=\frac{\ddot d/2+\dot q^2}{\ell},\qquad
 n_q=\frac{n_d/2+\dot q^2}{\ell},\\
 S_q=2q+\frac{x}{a_t}(1+a_t^{-2})\dot q-\ddot q/a_t^2,\\
 \rho=r\sqrt{(2+w^2/a_t^2)^2+(1+a_t^{-2})^2w^2x^2/a_t^2}.
\end{gathered}
\]

For \(0\le t\le3.6\), we use a uniform mesh and additional nodes for
interpolation. The samples are at \(t_i=i/40,\ 0\le i\le151\); a
subscript \(i\) denotes evaluation at that node. In the first table below,
the second column
encloses these values, with symmetric endpoints denoted \(\pm\).
The third column bounds
\(|F_{i+1}-2F_i+F_{i-1}|\) for \(0\le i\le150\), with the value
at \(-1\) supplied by parity. The fourth bounds
\(|F_{i+1}-F_i|\) for \(0\le i\le150\), when needed.

\begin{lemma}[Finite certificates]\label{lem:scalar-certificates}
The following sample bounds hold with exact rational endpoints.
\[
\begin{array}{c|c|c|c}
 &\text{value} & |\Delta^2| &|\Delta|\\ \hline
 d&[.3862,.5]&.000260&-\\
 \dot d&\pm.1338&.00092&.01037\\
 \ddot d&\pm.415&.00421&.03670\\
 \mathsf K&[-.00692,.02737]&.000148&-\\
 \dot{\mathsf K}&\pm.06542&.000673&-\\
 \ddot{\mathsf K}&\pm.2355&.00373&-\\
 n_{\mathsf K}&\pm.2355&.00402&-\\
 \mathsf C&[-.5,-.3413]&.00049&-\\
 \dot{\mathsf C}&\pm.2456&.00222&-\\
 \ddot{\mathsf C}&\pm.780&.0129&-\\
 n_{\mathsf C}&\pm.780&.0138&-\\
 R_v&\pm.318&.00167&-\\
 V_2&\pm1.06&.00186&-\\
 q&[.0777,.25525]&-&-\\
 \dot q&\pm.1908&-&.0127\\
 \ddot q&\pm.508&-&.047\\
 n_q&\pm.508&-&-
\end{array}
\]
The same nodes give the following simultaneous bounds:
\[
\begin{array}{c|cccc}
 i&\sup U_*&\inf(\lambda t_-^2+\mathsf Ct_-+\mathsf K)&\sup(\lambda t_+^2+\mathsf Ct_++\mathsf K)&
 \inf(2(-.37)\mathsf C-.37^2-4\lambda\mathsf K)\\\hline
 0{:}151&.537&.00090&-.00074&.132
\end{array}
\]
Here the entries for extrema are bounds in the indicated directions. For \(56\le i\le144\), sample bounds are \(0\le\mathsf K\le.019,\ |\dot{\mathsf K}|\le.0207,\ n_{\mathsf K}\ge.00315,\ -.5\le\mathsf C\le-.49512,\ |\dot{\mathsf C}|\le.0071,\ q\ge.2353\). For \(45\le i\le144\), \(q\ge.2204\).

For \(i=56,144\), respectively, the finite sums satisfy
\[
 x_i ((\mathsf K_i)_+ + .238/12800) -
 40\sum_{j=1}^{i} (a_{t_j}-a_{t_{j-1}})((\mathsf K_j)_+-(\mathsf K_{j-1})_+)
 \le .126,\ .335 .
\]
\end{lemma}

\begin{proof}
At each node, we first enclose \(x_i=s\sinh(i/40)\), \(a_{t_i}\),
\(\phi(x_i)\), and \(Y(x_i)\) by the rules below. These give
\(u,p,h,j\), and hence \(d,g',g'',\mathsf K,R_v,V_2\) by the stable
formulas. Next compute \(D_1,D_2,G_3,G_4\), and use \(P_F,Q_F\) to
obtain the dotted derivatives and \(n_F\). The enclosure \(d_i\le.5\)
at each node gives \(b-r^2-d_i\ge.0536\); only then do we evaluate
\(q_i\) and its derivative formulas. This is a nodewise check, not
yet a positivity assertion between nodes.

Taking interval hulls of these samples gives the value columns.
Subtraction gives the first and second differences on the index ranges
specified above. For the second difference at zero, use the odd
reflection for dotted first derivatives and \(R_v\), and the even
reflection for the other rows. Evaluating the four expressions in the
second table directly at the same nodes gives their simultaneous bounds.
The same arithmetic evaluates the two finite sums in the statement.
Interpolation will later turn them into integral bounds; here they
are only finite expressions in the enclosed node data. The primitive
enclosures are as follows.

\paragraph{Node coordinates and Gaussian density.}
Take
\(E=\big(\sum_{m=0}^{26}(i/160)^m/m!+[-10^{-27},10^{-27}]\big)^4\).
This encloses \(e^{i/40}\), so \(s(E-1/E)/2\) and \(s(E+1/E)/2\)
enclose \(x_i\) and \(a_{t_i}\), respectively; intersect the former
with \([0,90]\). For \(i\le73\), enclose \(\phi\) by the same
exponential sum with argument \(-x_i^2/128\) and the same error,
raised to the 64th power and divided by \(\sqrt{2\pi}\).
For \(i\ge74\), use \([0,10^{-26}]\). Indeed
\(x_{73}<11<x_{74}\), and \(e^{11^2/2}>10^{26}\), for instance from
\(e^{1/2}>1.645\). Every Taylor sum just used has a true argument
of absolute value at most \(1\). Thus
\(e^1/27!<3/27!<10^{-27}\) bounds its error before taking powers.

\paragraph{Gaussian tails.}
For \(i\le30\), take \(y_m=x_i(-x_i^2/2)^m/m!\) and
\(L=\sum_{m=0}^{39}y_m/(2m+1)+[0,y_{40}/81]\), then use
\(Y=(1/2-L/\sqrt{2\pi})/\phi\). This integrates the Taylor bounds
for a negative exponential. The terms can be formed recursively as
\(y_{m+1}=-y_m x_i^2/(2m+2)\).

For \(i\ge31\), division by the small Gaussian density is avoided.
Start with \(T_{34}=[0,34/x_i]\), put
\(T_m=m/(x_i+T_{m+1})\) for \(m=33,\ldots,1\), and use
\(Y=1/(x_i+T_1)\). To justify this enclosure, put
\(I_m(x)=\int_0^\infty y^m e^{-xy-y^2/2}\,dy\).
Integration by parts gives \(mI_{m-1}=xI_m+I_{m+1}\), so
\(T_m=I_m/I_{m-1}=m/(x+T_{m+1})\) and \(0<T_{34}<34/x\).
Also \(xI_0+I_1=1\), giving the formula for \(Y=I_0\).

\paragraph{Logarithmic profiles and square roots.}
The tail enclosure gives \(u=(Y\phi)\cap[0,1/2]\) and \(p=1-u\).
For \(h\), use
\(\sum_{m=0}^{44}u^m/(m+2)+[0,2u^{45}/47]\); the last interval
bounds the remaining positive series because \(u\le1/2\).
To evaluate the logarithms in \(j\), scale each argument by a power
of 2 into \([1,2.01]\). For a scaled argument \(z\), including
\(z=2\), take
\(2\sum_{m=0}^{24}y^{2m+1}/(2m+1)\), where \(y=(z-1)/(z+1)\),
with absolute error \(2|y|^{51}/(51(1-y^2))\). Correct for the
scaling by the corresponding multiple of \(\log2\).
Use \(\log\sqrt{2\pi}=\frac12\log(2\pi)\) and
\(3.14159265358979323846264338\le\pi\le3.14159265358979323846264339\).
These endpoints follow, for example, from the arctangent series at
\(1/5,1/239\) in \(\pi=16\arctan(1/5)-4\arctan(1/239)\).
For the square roots in the profiles, consecutive multiples of
\(10^{-26}\) enclosing each endpoint suffice.


Every addition, multiplication, reciprocal, and power is evaluated by
the full range of that operation on its interval arguments; endpoints
are rounded outwards. In particular the square of an interval crossing
zero has lower endpoint zero. The intersections and square-root rules
above preserve inclusion. They therefore give rational certificates,
without an assumption about the accuracy of numerical special functions.
For example, the two displayed finite integral sums lie respectively in
\([.1250504493,.1250504494]\) and
\([.3344848649,.3344848650]\), and the four simultaneous bounds have
strict slack before rounding to the displayed table.

The supplementary programs \texttt{scalar\_intervals.py} and
\texttt{scalar\_algebra.py} implement these rules with integer endpoints
on a \(10^{-70}\) grid and the specified \(10^{-26}\) square-root grid.
They also record each resulting interval. The formulas and remainder
bounds here specify the certificates independently of those programs.
A program-to-claim index and reproduction instructions for both
appendices are supplied in \texttt{verification/README.md}.
\end{proof}

\subsection{Analytic interpolation}

\begin{lemma}[A common analytic strip]\label{lem:scalar-strip}
As functions of \(t\), the profiles
\(d,\mathsf K,\mathsf C,R_v,V_2\) are analytic in a neighborhood of
\(\{t:|\Im t|\le .45\}\). On this strip, the first three have modulus
at most \(1300\), and the last two at most \(11000\).
\end{lemma}

\begin{proof}
First suppose \(\Re t\ge0\), and write \(x=A+iy=s\sinh t\).
Then \(A\ge0\), and the elementary bounds
\(\sin(.45)<.435\), \(\tan(.45)<.5\) give
\[
 y^2\le2.5+.25A^2.
\]
The Laplace integrals for \(Y\) and \(B=-Y'\) give
\[
 |Y|\le Y(A)\le\min(1.254,1/A),\qquad
 |B|\le\min(1,1/A^2),
\]
where expressions containing \(1/A\) are omitted at \(A=0\).
The corresponding tail integral for \(u=1-p\) gives
\[
 |u|\le e^{y^2/2}u(A)\le1.75,
\]
since \(u(A)e^{A^2/8}\) decreases on \([0,\infty)\).
Indeed \(Au(A)\le\phi(A)\) follows by integrating the Gaussian tail,
and differentiating the product proves this monotonicity.

If \(A\le.9\), then \(A|y|<\pi/2\). Horizontal integration of
\(\phi\) from \(iy\) to \(A+iy\), starting with
\(\Re p(iy)=1/2\), gives \(\Re p\ge1/2\).
If \(A\ge.9\), then
\[
 |u|\le u(.9)e^{1.25+.125\cdot.9^2}<.80,
 \qquad
 u(.9)\le\tfrac12-\frac{.9-.9^3/6}{\sqrt{2\pi}}.
\]
Thus \(\Re p>0\) throughout the right-half image.
The identity
\[
 h(u)=\int_0^1\frac{v}{1-uv}\,dv
\]
shows analyticity and \(|h|\le2.5\): the denominator has real part at
least \(1/2\) in the first region and modulus at least \(.2\) in the
second.

We also need a nonvanishing bound for \(Y\). Its Laplace transform and
the characteristic function \(e^{-A|s'|}\) of a Cauchy variable \(T_A\)
of scale \(A\) give
\[
 \Re Y(A+iy)=\sqrt{\pi/2}\,
      \mathbb E e^{-(y+T_A)^2/2}>0.
\]
At \(A=0\), take \(T_A=0\). If \(A\le2\), with probability at least
\(1/4\) one has
\(|y+T_A|\le\max(|y|,A)\le2\). For instance, when \(y\ge0\)
the permitted interval contains \([-A,0]\); reflection handles the
other sign. Hence \(\Re Y\ge\sqrt{\pi/2}e^{-2}/4>.02\).
If \(A\ge2\), the mean of the integration variable under the
Laplace weight is at most \(1/A\), because the Gaussian factor is a
decreasing tilt of an exponential density. Consequently
\[
 \left|\frac{Y(A+iy)}{Y(A)}-1\right|\le\frac{|y|}{A}
 \le\sqrt{7/8}<.936,
 \qquad
 Y(A)\ge\frac{e^{-1-1/(2A^2)}}{A}>\frac{.324}{A}.
\]
In both regions,
\[
 |Y|\ge\frac{.02}{1+A}.
\]

The principal logarithms of \(p\) and \(Y\) are therefore analytic
here. The analytic logarithm of \(1-p=\phi Y\) is
\(\log Y-x^2/2-\log\sqrt{2\pi}\), which agrees with its real value
on the positive axis. This specifies the branch in \(j\).
The preceding bounds yield
\[
 |j|\le .625A^2+\log(1+A)+10.5,\qquad
 |Bj|\le12,\qquad |Y^2j|\le19.
\]
For example, use \(|\log Y|\le\log50+\log(1+A)+\pi/2\),
\(|p|\le2.75\), and then split at \(A=1\) for the last two bounds.
The stable formulas, together with \(|1-pB|\le3.75\) and
\(|xY|=|1-B|\le2\), now imply
\[
 |d|<48,\qquad |\mathsf K|<99,\qquad |xg'|<76,
 \qquad |\mathsf C|\le48+6.5(76+99)<1300.
\]
Also
\[
 |\phi/p|\le7e^{-.375A^2},\qquad
 |B/Y|\le100,\qquad |B/Y^2|\le10000.
\]
Thus \(|R_v|\le107\). Since
\(|x|\le1.25A+1.6\) and \(Ae^{-.375A^2}\le1\),
\(|x\phi/p|<20\), giving \(|V_2|\le2+20+49+10000<11000\).

Reflection gives the left-half bounds. On the imaginary segment in
question, both \(p\) and \(1-p\) have real part \(1/2\), and
\(\Re Y>0\). The formulas are therefore regular on a neighborhood of
that segment. The right-half expressions and their parity reflections
agree near zero and hence by analytic continuation across the segment.
This proves the asserted full-strip analyticity and bounds.
\end{proof}

\begin{lemma}[Interpolation from second differences]
\label{lem:scalar-interpolation}
Put \(\mathcal R(M)=.536M+.00001\). For any of
\[
 F,\dot F,\ddot F,n_F\quad(F=d,\mathsf K,\mathsf C),
 \qquad R_v,V_2,
\]
the error from its central chord, the affine interpolant through the
two endpoints of a mesh cell in \([0,3.6]\), is
at most \(\mathcal R(M)\), where \(M\) bounds its absolute second
differences whose three nodes lie in the sixteen-node stencil for that cell.
\end{lemma}

\begin{proof}
Fix one of the functions in the statement and denote it by \(\Psi\).
Use the 16 nodes \(i-7,\ldots,i+8\) for the cell \([t_i,t_{i+1}]\),
reflecting by parity when necessary. Let \(\mathcal I_8\) be its
interpolating polynomial on this stencil, of degree at most 15.
We first bound \(\Psi-\mathcal I_8\), and then compare
\(\mathcal I_8\) with the central chord. The first error is at most
\[
 4\cdot10^5 \prod_{j=0}^7((j+1/2)/(40\cdot .36))^2 < .00001
\]
in the cell. Indeed Lemma~\ref{lem:scalar-strip} and Cauchy's estimates
from the strip of half-width \(.45\) to that of half-width \(.36\)
bound, across the gap \(.09\),
\(\dot F\) by \(1300/.09\), \(\ddot F\) by
\(2\cdot1300/.09^2\), and \(n_F\) by their sum, which is less than
\(4\cdot10^5\). Here \(|\tanh t|\le1\) on \(|\Im t|\le.36\).
Apply the real interpolation remainder and Cauchy's estimate of order
16 on the inner strip. The paired nodal factors attain their maximum
at the midpoint of the cell; the displayed bound is less than
\(.000007336661\).

For the comparison with the chord, normalize the cell by
\(t=t_i+\delta\theta\), where \(\delta=1/40\) and \(0\le\theta\le1\),
and put \(\Psi_k=\Psi(t_i+k\delta)\).
For \(1\le l\le8\), let \(\mathcal I_l\) interpolate these values at
\(k=1-l,\ldots,l\); thus \(\mathcal I_1\) is the central chord.
For \(2\le l\le8\), let \(\mathcal J_l^-\) omit the rightmost node \(l\), and let
\(\mathcal J_l^+\) omit the leftmost node \(1-l\).
The interpolation identity
\[
 \mathcal I_l(\theta)
 =\frac{l-\theta}{2l-1}\mathcal J_l^-(\theta)
  +\frac{\theta+l-1}{2l-1}\mathcal J_l^+(\theta)
\]
follows by checking the nodal values and the degree. Its weights are
nonnegative and sum to one on the central cell.

The common nodes of \(\mathcal J_l^-\) and \(\mathcal J_l^+\) are
exactly those of \(\mathcal I_{l-1}\). With
\(\Delta\Psi_k=\Psi_{k+1}-\Psi_k\), Newton interpolation therefore gives
\[
\begin{aligned}
 \mathcal J_l^--\mathcal I_{l-1}
  &=\frac{\Delta^{2l-2}\Psi_{1-l}}{(2l-2)!}
       \prod_{k=2-l}^{l-1}(\theta-k),\\
 \mathcal J_l^+-\mathcal I_{l-1}
  &=\frac{\Delta^{2l-2}\Psi_{2-l}}{(2l-2)!}
       \prod_{k=2-l}^{l-1}(\theta-k).
\end{aligned}
\]
There is no mesh factor in these formulas because the nodes in the
\(\theta\) coordinate have spacing one. Equivalently, the powers of
\(\delta\) in the nodal product cancel those in the divided difference.
Expanding \(\Delta^{2l-2}=\Delta^{2l-4}\Delta^2\) and using the
second-difference bound gives
\[
 |\Delta^{2l-2}\Psi_k|\le 2^{2l-4}M.
\]
Pair the nodal factors about \(1/2\). On \(0\le\theta\le1\),
\[
 \left|\prod_{k=2-l}^{l-1}(\theta-k)\right|
 =\prod_{j=0}^{l-2}\big((j+\tfrac12)^2-(\theta-\tfrac12)^2\big)
 \le\prod_{j=0}^{l-2}(j+\tfrac12)^2.
\]
The convex-blend identity bounds \(|\mathcal I_l-\mathcal I_{l-1}|\)
by the product of these two bounds divided by \((2l-2)!\).
Summing over the seven successive pairs of added nodes yields
\[
 |\mathcal I_8-\mathcal I_1|
 \le M\sum_{l=2}^8
      \frac{2^{2l-4}}{(2l-2)!}\prod_{j=0}^{l-2}(j+1/2)^2
 =\frac{4387}{8192}M<.536M.
\]
Adding the analytic remainder proves the claimed chord error.
For the final cell \(i=143\), the largest stencil index is \(151\);
near zero, parity supplies the negative indices. These stencils cover
the entire interval \([0,3.6]\).
\end{proof}

\subsection{Bounds between the nodes}

Apply Lemma~\ref{lem:scalar-interpolation} using the third column of
Lemma~\ref{lem:scalar-certificates} on its first thirteen lines. In particular \(d\le.50015\), so \(b-r^2-d\ge.05345>0\) on this finite region. This conclusion uses only the \(d\) table and its interpolation bound, and therefore precedes all estimates involving \(q\) or \(1/\ell\).

We now estimate the chord errors for \(q\) and its derivatives.
Set \(\delta=1/40\) and \(l=1/(2\ell)\); the positive radicand just
proved licenses this division throughout the finite region.
The interpolated bounds for \(d,\dot d,\ddot d\), together with
\[
 \dot l=2l^3\dot d,\qquad
 \ddot l=2l^3\ddot d+12l^5\dot d^2,
\]
give \(l\le2.164\), \(|\dot l|\le2.8\), and
\(|\ddot l|\le19.2\). The chord error for \(l\) is consequently at
most \(e=19.2\delta^2/8=19.2/12800\).

For the products below, let \(J_E,J_F,J_{EF}\) denote the affine
chords of \(E,F,EF\) on a cell, and put
\(\theta=(t-t_i)/\delta\). The exact identity
\[
 EF-J_{EF}
 =E(F-J_F)+J_F(E-J_E)
  -\theta(1-\theta)(E_{i+1}-E_i)(F_{i+1}-F_i)
\]
shows that errors \(\epsilon_E,\epsilon_F\) for the two factors give
the product error
\[
 \sup|E|\,\epsilon_F+
 \max(|F_i|,|F_{i+1}|)\,\epsilon_E+
 \tfrac14|E_{i+1}-E_i|\,|F_{i+1}-F_i|.
\]
Apply this first to \(\dot q=l\dot d\), using the first-difference
column of the sample table and \(|l_{i+1}-l_i|\le2.8\delta\).
It gives
\[
 2.164{\cal R}(.00092)+.1338e+
       (2.8\delta)\cdot.01037/4<.0016.
\]
Since the endpoint values of \(\dot q\) have modulus at most \(.1908\)
and first differences at most \(.0127\), the corresponding error
for \(\dot q^2\) is at most
\[
 (2\cdot.1908+.0016)\cdot.0016+.0127^2/4<.000654.
\]
We may therefore apply the product bound to
\(\ddot q=l\ddot d+2l\dot q^2\). The resulting error is bounded by
\[
 \begin{aligned}
 &2.164{\cal R}(.00421)+.415 e+(2.8\delta)\cdot.03670/4\\
 &\qquad{}+2\big[2.164\cdot.000654+.1908^2 e
       +(2.8\delta)(2\cdot.1908\cdot.0127)/4\big]<.0095,
 \end{aligned}
\]
where the first three terms bound the contribution of \(l\ddot d\),
and twice the bracketed expression bounds that of \(2l\dot q^2\).
Combining this error with the sample bound
\(|\ddot q_i|\le.508\) gives a chord error
\((.508+.0095)\delta^2/8<.000042\) for \(q\).

For \(n_q=\ddot q-\tanh(t)\dot q\), use the product bound once more.
The first and second derivatives of \(\tanh t\) have modulus at most
\(1,1\), respectively, so the error is less than
\[
 .0095+.0016+.1908/12800+\delta\cdot.0127/4<.012.
\]
For \(S_q\), use
\(S_q=2q+\tanh(t)(1+a_t^{-2})\dot q-a_t^{-2}\ddot q\).
The first and second derivative bounds for \(a_t^{-2}\) are
\(1/s^2,2/s^2\), and those for
\(\tanh(t)(1+a_t^{-2})\) are \(1.08,2\).
The three terms in this expression therefore have total error at most
\[
 \begin{aligned}
 &2\cdot.000042+(1+1/s^2)\cdot.0016+
 .1908\cdot2/12800+(1.08\delta)\cdot.0127/4\\
 &\qquad{}+\frac{.0095+.508\cdot2/12800+\delta\cdot.047/4}{s^2}<.003.
 \end{aligned}
\]
We have obtained the five errors
\(.000042,.0016,.0095,.012,.003\) for
\(q,\dot q,\ddot q,n_q,S_q\).

We can sharpen value errors of \(\mathsf K,\mathsf C\) to \(.238\delta^2/8,.79\delta^2/8\) using the established derivative bounds.

The expression \(U_*\) is a convex function of the three arguments
\((-n_{\mathsf K},\dot{\mathsf K},\dot{\mathsf C})\): its first term
is a positive part and the others are positive multiples of squares
of linear forms. On the chord joining two sample triples it is thus
at most the larger endpoint value, bounded by \(.537\).
It remains to account for the errors in those three arguments.
Put \(\epsilon_c={\cal R}(.00222)\) and
\(E_*={\cal R}(.000673)+t_c\epsilon_c\).
The positive-part term changes by at most
\(1.507{\cal R}(.00402)\); applying
\(|a^2-b^2|\le|a-b|(2|b|+|a-b|)\) to each square gives the
following total addition to \(.537\):
\[
 1.507{\cal R}(.00402)+\frac{12\cdot .80}{m_*^2}
 \left(1.25 E_*\big[2(.06542+t_c\cdot.2456)+E_*\big]
             +5t_r^2\epsilon_c(2\cdot .2456+\epsilon_c)\right)<.010 .
\]
Consequently \(U_*<.547<.54r^2w^2\) on the finite region.

As for the three inequalities sampled together with \(U_*\), upper absolute errors from the chords are \(<.000024,.000024,.0001\), respectively, using \(.238/12800<.000019,\ .79/12800<.000062\) for the \(\mathsf K,\mathsf C\) errors. Hence the margins in that table suffice. These estimates prove the finite-range part of Lemma~\ref{lem:scalar-ranges}.
For the restricted range \(t\ge1.4\), the central sample values obey
\(n_{\mathsf K}\ge.00315>\mathcal R(.00402)\), so this sign persists
between nodes. The condition \(x\ge5\) is covered because
\(s\sinh(45/40)<5\).

To interpret the two finite integral sums, let \(f_i=(\mathsf K_i)_+\).
Convexity of positive part bounds the positive part of a chord of
\(\mathsf K\) by the chord through \(f_i\). Add the established
value error \(.238/12800\) and integrate against \(dx=s\cosh t\,dt\).
Summation by parts gives precisely the sums in
Lemma~\ref{lem:scalar-certificates}: the endpoint term is \(x_if_i\),
and the contribution of the slope \(40(f_j-f_{j-1})\) on the
\(j\)th cell is
\(-40(f_j-f_{j-1})(a_{t_j}-a_{t_{j-1}})\), because
\(\int x\,dt=s\cosh t\). The added constant error contributes
\(x_i\cdot.238/12800\). Their upper bounds are therefore
valid for the full integrals, not merely for a quadrature at the nodes.
It remains to control the infinite tail and the layer inequalities.

\subsection{Uniform estimates on the tails}

\begin{lemma}[Tail derivative bounds]\label{lem:scalar-tail}
For \(x\ge64\), set \(D=x\partial_x\), \(\xi=x^{-2}\),
and \(L=\log x+\frac12\log(2\pi)-\frac12\).
For each row function \(F\), the entries are upper bounds for
\(|D^jF|/\xi\), uniformly on this range:
\[
\begin{array}{c|ccc}
 &j=0&j=1&j=2\\\hline
 \mathsf K-\xi(L-\tfrac32)&.0064&.027&.11\\
 d-\tfrac12&.51&1.02&2.1\\
 \mathsf C+\tfrac12&.55&1.16&2.61
\end{array}
\]
These estimates supply the tail part of the proof of
Lemma~\ref{lem:scalar-ranges}. On \(x\ge64\) we also have
\(q\in[.255,.2556]\), \(S_q\in[.509,.512]\),
\(V_2\in[.999,1.001]\).
\end{lemma}

\begin{proof}
It suffices to work on \(x\ge64\): the endpoint of the finite region
is \(x_{144}>65.8>64\), so the two arguments overlap. Put \(D=x\partial_x,\ \xi=x^{-2},\ \xi_0=1/4096,\ m=xY=1-\xi n\), with
\[
 n=\int_0^\infty y e^{-y-\xi y^2/2}\,dy,\qquad L=\log x+\tfrac12\log(2\pi)-\tfrac12,\qquad
 Z=\log x+\tfrac12\log(2\pi)-\log m-p+p^2h.
\]
Useful exact expressions are
\[
\begin{aligned}
 d&=h+\xi(\tfrac12 n^2-2p h n)+\xi^2 n^2 Z,\qquad
 g=\tfrac14+\tfrac14(1-m^2)-\tfrac12 \xi m^2 Z,\\
 \mathsf K&=(h-\tfrac12)+\xi\big(\mathcal A+(m^2+\xi n^2)Z\big),\qquad
 \mathcal A=\tfrac12 n^2-2p h n-n+\tfrac12\xi n^2 .
\end{aligned}
\]
We prove the three rows by keeping each derivative bound and its
prefactor explicit. For a nonnegative numerical vector
\(V=(V_j)_{j=0}^4\), write
\[
 F\preccurlyeq aV
 \quad\hbox{to mean}\quad
 |D^jF(x)|\le a(x)V_j
 \quad(0\le j\le4,\ x\ge64).
\]
The prefactor \(a\) is a pointwise bound; the notation does not
differentiate \(a\). If \(F\preccurlyeq aA\) and
\(G\preccurlyeq bB\), the product rule gives
\(FG\preccurlyeq ab(AB)\), where juxtaposition now denotes
binomial convolution:
\[
 (AB)_j=\sum_{i=0}^j\binom ji A_iB_{j-i}.
\]
Use the numerical vectors
\(P_j=2^j\), \(I_j={\bf1}_{j=0}\), and \(J_j={\bf1}_{j=1}\).
In particular \(\xi\preccurlyeq\xi P\), since \(D\xi=-2\xi\).
Thus multiplying a function by \(\xi\) introduces both that prefactor
and a convolution by \(P\). All vector inequalities below are
componentwise.

\paragraph{The bounds for \(n\) and \(m\).}
Differentiating the integral for \(n\), with
\(D=-2\xi\partial_\xi\), gives
\(|n^{(i)}_\xi|\le(2i+1)!/2^i\) and \(|1-n|\le3\xi\).
For \(1\le j\le4\), expand \((\xi\partial_\xi)^j\).
After division by \(2^j\xi\), each resulting bound is at most
\[
 3+7\cdot30\xi_0+6\cdot630\xi_0^2+22680\xi_0^3<3.1.
\]
Consequently \(n-1\preccurlyeq3.1\xi P\). Set
\(N=I+3.1\xi_0P\), so \(n\preccurlyeq N\).
The identity \(m-1=-\xi n\) then gives
\[
 m-1\preccurlyeq\xi PN,\qquad PN\le1.02P.
\]
With \(M=I+1.02\xi_0P\), we also have \(m\preccurlyeq M\).

\paragraph{The logarithmic term.}
The series for \(-\log m\) is a series in \(1-m\).
For a power of order \(k\), convolution satisfies
\((P^k)_j=k^jP_j\). Including the factor \(1/k\) from the logarithm
and using \(j\le4\), its derivative bounds are therefore controlled by
\[
 -\log m\preccurlyeq
 \xi\cdot1.02P\sum_{k\ge1}k^3(1.02\xi_0)^{k-1}.
\]
The vector on the right is at most \(1.04\xi P\), using
\(\sum_{k\ge1}k^3z^{k-1}=(1+4z+z^2)/(1-z)^4\);
the series and its derivatives converge absolutely in this tail.

The remaining corrections in \(Z\) are Gaussian. The product rule
applied to \(u=\phi m/x\) gives
\(|D^ju|\le10^{-18}\xi^2\) for \(j\le4\), using
\(x^{20}e^{-x^2/2}<10^{-200}\). Substitution into the series for \(h\)
gives
\[
 h-\tfrac12,\quad ph-\tfrac12,\quad p-1
       \preccurlyeq10^{-6}\xi^2P.
\]
Since \(Z-L=-\log m-p+p^2h+\tfrac12\), these bounds imply
\(Z-L\preccurlyeq1.05\xi P\).
We also need a bound for \(Z\) itself, whose logarithmic growth
prevents a uniform absolute bound on the tail. Put
\(H_0=4.59I+J+1.05\xi_0P\). Since
\(DL=1\), \(D^jL=0\) for \(j\ge2\), and \(\xi L\) decreases on
\(x\ge64\), the precise bounds represented by these vectors are
\[
 |D^jn|\le N_j,\qquad |D^jm|\le M_j,\qquad
 \xi|D^jZ|\le\xi_0(H_0)_j,\qquad 0\le j\le4.
\]
Thus \(Z\preccurlyeq(\xi_0/\xi)H_0\). The last prefactor will cancel
one power of \(\xi\) when we bound the exact formulas above.

\paragraph{The algebraic corrections.}
The two corrections needed for \(\mathsf K\) are
\(\mathcal A+1.5\) and \(m^2+\xi n^2-1\).
To bound the first, use the identity
\[
 \mathcal A+1.5
 =\tfrac12(n-1)(n+1)-2(n-1)+\tfrac12\xi n^2
       -2(ph-\tfrac12)n.
\]
For the second, write \(m^2-1=(m-1)(m+1)\).
The product rule then bounds their derivative vectors, after division
by \(\xi\), by the first and second lines below, respectively:
\[
\begin{aligned}
 &1.55P(N+I)+6.2P+.5P N^2+2\cdot10^{-6}\xi_0 P N \ \le\ 10P,\\
 &1.02P(M+I)+P N^2\ \le\ 3.2P .
\end{aligned}
\]
The comparisons on the right hold through order 4, as follows by
expanding \(N,M\) and using \((P^k)_j=k^jP_j\).
The bound \(10\xi P\) also applies to
\(\tfrac12n^2-2phn+\tfrac12\). Its expansion replaces
\(-2(n-1)\) by \(-(n-1)\) and omits \(\tfrac12\xi n^2\), so the
same nonnegative bound suffices.

\paragraph{The remainder in \(\mathsf K\).}
Subtract the main term from its exact expression:
\[
 \mathsf K-\xi(L-\tfrac32)
 =(h-\tfrac12)+\xi\big[
   (\mathcal A+\tfrac32)+(Z-L)+(m^2+\xi n^2-1)Z\big].
\]
The first two terms inside brackets have combined bound
\(11.05\xi P\); the last has bound \(3.2\xi_0PH_0\).
Multiplication by the outside factor \(\xi\), followed by
\(\xi\le\xi_0\), therefore bounds the derivative vector of this
remainder by \(\xi\) times
\[
 \xi_0 (10^{-6}P + P(11.05P+3.2P H_0)).
\]
Its first three components are at most \((.0064,.027,.11)\),
as in the first table row.

\paragraph{The remainder in \(d\).}
Isolate \(-\xi/2\) in the formula for \(d-\tfrac12\):
\[
 d-\tfrac12
 =(h-\tfrac12)-\tfrac12\xi+
    \xi(\tfrac12n^2-2phn+\tfrac12)+\xi^2n^2Z.
\]
The four terms give \(d-\tfrac12\preccurlyeq\xi E_d\), where
\[
 E_d=\tfrac12P+
       \xi_0(10^{-6}P+10P^2+P^2N^2H_0).
\]
The last term uses \(\xi^2\preccurlyeq\xi^2P^2\) and
\(Z\preccurlyeq(\xi_0/\xi)H_0\).
The first three components of \(E_d\) are at most
\((.51,1.02,2.1)\), proving the second row.

\paragraph{From \(g\) to the remainder in \(\mathsf C\).}
Using \(m=1-\xi n\), rewrite
\[
 g-\tfrac14=\tfrac12\xi n-\tfrac14\xi^2n^2-\tfrac12\xi m^2Z.
\]
The product bounds give the absolute estimate
\(g-\tfrac14\preccurlyeq\xi_0E_g\), with the constant prefactor
\(\xi_0\), where
\[
 \xi_0E_g
 =\xi_0P(\tfrac12N+.25\xi_0PN^2+.5M^2H_0).
\]
This estimate is not in units of the varying \(\xi\).
Now \(g''=\xi(D^2-D)g\), so two more \(D\)-derivatives and one
factor \(\xi\) are required. Explicitly, for \(0\le j\le2\),
\[
 |D^jg''|
 \le \xi\xi_0\sum_{i=0}^j\binom ji P_i
       \big((E_g)_{j-i+2}+(E_g)_{j-i+1}\big).
\]
Since \(\mathsf C+\tfrac12=-(d-\tfrac12)+6.5g''\), its derivative
bounds in units of \(\xi\), through order 2, are
\[
 E_d+6.5\xi_0P
       \big((E_g)_{j+1}+(E_g)_{j+2}\big)_{j=0}^2.
\]
Their components are at most \((.55,1.16,2.61)\), proving the
third row. This completes the derivative table.

\paragraph{Conversion to the real derivative bounds.}
The \(D\)-bounds give the derivatives used in
Lemma~\ref{lem:scalar-ranges} through
\[
 \dot F=(a_t/x)DF,\qquad
 n_F=(a_t/x)^2(D^2-D)F,\qquad
 \ddot F=n_F+DF,\qquad a_t/x\le1.002.
\]
In particular,
\((D^2-D)(\xi(L-1.5))=\xi(6(L-1.5)-5)\).
The first row of the table therefore gives
\[
 (D^2-D)\mathsf K
 \ge\xi\big(6(L-1.5)-5-.11-.027\big)>0,
\]
because \(L\ge L(64)>4.5\). Hence \(n_{\mathsf K}>0\).
For upper bounds, use \(L(64)<4.59\) and the fact that
\(\xi(L-1.5)\) and \(\xi(2L-4)\) decrease on this tail. The table
and the conversion formulas give
\[
\begin{aligned}
 0&\le\mathsf K\le\xi_0(4.59-1.5+.0064),\quad
 |\dot{\mathsf K}|\le1.002\xi_0(2\cdot4.59-4+.027),\quad
 |\ddot{\mathsf K}|<.02,\\
 |\mathsf C+.5|&\le .55\xi_0,\quad |\dot{\mathsf C}|\le1.002\cdot1.16\xi_0,\quad
 |n_{\mathsf C}|,\ |\ddot{\mathsf C}|<.002.
\end{aligned}
\]
These imply the global and restricted derivative bounds for
\(\mathsf K,\mathsf C\) on the tail.

\paragraph{The \(q\) and layer-profile ranges.}
The second row of the table first gives
\[
 \ell^2=b-r^2-d\ge.0536-.51\xi_0>.231^2.
\]
This tail bound on the radicand is independent of the estimates for
\(q\), and permits division by \(\ell\) in its derivative formulas.
Substitution of the value and derivative bounds for \(d\) gives
\[
 k-\sqrt{.0536+.51\xi_0}\le q\le k-\sqrt{.0536-.51\xi_0},\qquad
 |\dot q|<.00055,\quad |n_q|<.0017.
\]
Thus \(q\in[.255,.2556]\). The identity
\(S_q=2q+\tanh(t)\dot q-a_t^{-2}n_q\) then yields
\(S_q\in[.509,.512]\). The bounds for \(\dot q,n_q\) also give
the required bound for \(\ddot q=n_q+\tanh(t)\dot q\).

For the two derivatives of \(v\), the integral for \(n\) gives
\(1-3\xi\le n\le1\), while \(m=1-\xi n\) gives
\(1-\xi\le m\le1\). Consequently
\[
 V_2=2-(\phi/p)(x+\phi/p)-n/m^2\in[.999,1.001],
 \qquad R_v=\phi/p-n/(xm).
\]
The Gaussian term in \(V_2\) is smaller than \(.0001\).
The same Gaussian estimate and \(m\ge1-\xi_0\) bound \(|R_v|\)
by \(.319\), as required.

\paragraph{The covariance inequalities and \(U_*\).}
The value estimates just obtained imply
\(0\le\mathsf K\le.0008\) and
\(-.501\le\mathsf C\le-.499\).
Inserting these intervals in
\(\lambda t_\pm^2+t_\pm\mathsf C+\mathsf K\) gives the lower bound
\(.010\) for \(t_-\) and the upper bound \(-.027\) for \(t_+\).
For the third covariance inequality, the same intervals give
\[
 2(-.37)\mathsf C-.37^2-4\lambda\mathsf K
 \ge.74\cdot.499-.37^2-2.6\cdot.0008>m_*^2.
\]
Finally \(n_{\mathsf K}>0\) removes the positive-part term in \(U_*\).
Using \(|\dot{\mathsf K}|<.002\) and
\(|\dot{\mathsf C}|<.001\) in its two squares gives \(U_*<.001\).

\paragraph{The positive-part integral.}
The first table row and
\(\log\sqrt{2\pi}-2+.0064<0\) give
\(\mathsf K\le(\log x)x^{-2}\). Hence
\[
 \int_{x_{144}}^\infty\mathsf K_+
 \le\int_{64}^\infty\frac{\log x}{x^2}\,dx
 =\frac{\log64+1}{64}<.086.
\]
Combining this with the finite upper bound \(.335\) gives a total
less than \(.422\). All derivatives needed in the finite argument
are also bounded by the displayed tail estimates. Reflection completes
the proof of Lemma~\ref{lem:scalar-ranges}.
\end{proof}

For every integer \(j\ge0\) and \(x\ge1\), differentiating the
Laplace integral gives
\[
 Y^{(j)}(x)=(-1)^j\int_0^\infty y^j e^{-xy-y^2/2}\,dy,
 \qquad |Y^{(j)}(x)|\le j!x^{-j-1}.
\]
Together with \(Y(x)\asymp x^{-1}\) on this half-line, the stable
formulas, reflection, and the uniform lower bound for the radicand show
that \(d,g,v,\mathsf K,\mathsf C,\pi,q\) are smooth and have derivatives
of at most polynomial growth. These global growth assertions concern
the profiles, not the temporary factors \(X,Y\) separately.


\subsection{The pointwise layer inequalities}

\begin{proof}[Proof of Lemma~\ref{lem:scalar-layer}]
We apply the real bounds separately to
\(t_x=\operatorname{arsinh}(x/s)\) and
\(t_z=\operatorname{arsinh}(z/s)\).
The following table encloses the values of the even profiles
\(q,S_q,V_2,\rho\), not merely their samples. Each row specifies an
interval of \(|t|\) in units of \(1/40\), and every displayed bound
is multiplied by 1000.

For \(q,S_q,V_2\), take the sample minima and maxima on each closed
row through index \(144\), and enlarge by the errors
\(.000042,.003,.00101\), respectively. The last row also includes
the tail bounds of Lemma~\ref{lem:scalar-tail}.
For \(\rho\), the exact formula makes an endpoint bound possible.
Put \(y=\tanh^2t\) and \(v_1=(1-y)/s^2\). Differentiating gives
\[
 \frac{d}{dy}\frac{\rho^2}{r^2w^2}
 =(1+v_1)^2-\frac{2(2+w^2v_1)}{s^2}
             -\frac{2y(1+v_1)}{s^2}
 \ge1-\frac6{s^2}-\frac{2(w^2+1)}{s^4}>.27.
\]
Here \(0\le y\le1\) and \(0\le v_1\le1/s^2\).
Thus \(\rho\) increases for \(t\ge0\), giving its upper endpoint
bound in each row, including the limiting bound in the last row.
\[
\begin{array}{c|cc cc cc c}
 & q_{\min}&q_{\max}& S_{q,\min}& S_{q,\max}& V_{2,\min}& V_{2,\max}& \rho_{\max}\\\hline
 0{:}4&77&81&113&132&725&743&804\\
 4{:}8&80&88&125&166&740&785&814\\
 8{:}12&87&100&159&220&782&848&831\\
 12{:}16&99&116&213&288&845&920&853\\
 16{:}20&115&134&281&359&917&986&877\\
 20{:}28&133&171&352&464&983&1055&927\\
 28{:}40&170&211&457&504&1041&1059&992\\
 40{:}68&210&245&497&513&1009&1044&1072\\
 68{+}&244&256&506&514&999&1012&1104
\end{array}
\]
In particular, the table proves \(.113\le S_q\le.514\) and
\(\rho\le1.104\). Number its rows from 1 to 9, and write
\(q_i^\pm,S_i^\pm,V_i^\pm,\rho_i^+\) for their endpoints after
division by 1000. Let \(i_x,i_z\) be the rows containing
\(|t_x|,|t_z|\); these two row choices are independent.

We first enclose \(H\) on this pair of rows.
Its scalar part \(4c+4kq+2(k-q)S\) increases in both \(q\) and \(S\),
because its partial derivatives are \(4k-2S>0\) and \(2(k-q)>0\).
The inequalities \(q<k\) and \(S_q<2k\) used here follow from the
table. Also \(|\pi(x)|=r\), so
\(|2S_\pi(z)\cdot\pi(x)|\le2r\rho_{i_z}^+\).
Consequently \(H\in[H_-,H_+]\), where
\[
\begin{aligned}
 H_-&=4c+4kq_{i_x}^-+
        2(k-q_{i_x}^-)S_{i_z}^- -2r\rho_{i_z}^+,\\
 H_+&=4c+4kq_{i_x}^++
        2(k-q_{i_x}^+)S_{i_z}^+ +2r\rho_{i_z}^+.
\end{aligned}
\]
At the same time \(V_2(x)\in[V_{i_x}^-,V_{i_x}^+]\).

To verify the first layer inequality, subtract its right side from
its left side and use \(J=H+V_2(x)-2\kappa\).
Since \(S_q(x),S_q(z)\) are positive, replacing their product by
\(S_{i_x}^+S_{i_z}^+\) gives the lower bound
\[
\begin{aligned}
 \mathscr D_{i_x,i_z}(H,V)
 ={}&H+t_c(H+V-2\kappa)
       -\left(\frac{t_r}{2h_s}+\frac1\alpha\right)
          (H+V-2\kappa)^2\\
 &-2(b+\eta+(b+\eta-\kappa)t_-)
   -2b_mS_{i_z}^+S_{i_x}^+
   -2e_0(2t_r)^2-\frac{t_rh_s}{2}.
\end{aligned}
\]
This is a concave function of the two coordinates \(H,V\).
Every point of
\([H_-,H_+]\times[V_{i_x}^-,V_{i_x}^+]\) is a convex combination
of its four vertices, so its value is at least the smallest vertex
value. There are \(9\cdot9=81\) pairs of rows and four vertex
evaluations for each pair. The resulting lower bounds, multiplied
by 1000 and minimized over all rows of \(x\), are
\[
 11,\ 14,\ 33,\ 64,\ 63,\ 39,\ 54,\ 45,\ 40
\]
in the order of the rows of \(z\). Each is positive, proving the
strict inequality on every rectangle.

For the last inequality in the lemma,
\(4c+2k(S_q(x)+S_q(z))-S_q(x)S_q(z)\) increases in each \(S_q\)
argument, again because the other argument is less than \(2k\).
The vector product is at most \(\rho(x)\rho(z)\le1.104^2\).
Thus its left side is bounded below by
\[
 4c+4k\cdot.113-.113^2-1.104^2>.25.
\]


Finally the auxiliary estimates follow directly from rational
arithmetic and the logarithm enclosure in
Lemma~\ref{lem:scalar-certificates}:
\[
 b+\eta+\kappa t_-+.125>0,\qquad
 \frac{1.053\kappa^2}{8(b+\eta+\kappa t_-+.125)}<.256,
 \qquad .256+\frac{(\log2+.25)2.26}{4}<.80.
\]
All rectangle margins are strictly positive. This proves the pointwise
assertions for every pair of real arguments and completes the scalar
estimates.
\end{proof}
